\subsection{The terminal pairs $(2,b)$ for every $b\ge4$ at a unit root}\label{proof:R2-C035}
For every $b\ge4$ and $L\ge1$, the fixed actual direction
\[
H_D=E_rpT^L(E_th_2f_b),\qquad U_D=T^2H_D
\]
at a unit internal root and an outer path of exactly two edges satisfies $I_q(U_D)^3\le e^{-d/500}$. Here $E_t,E_r$ are the genuine equal-pair products at the indicated original vertices. The direction is the same on every analytic subdomain below. For $L\ge2$, the ordinary fourth-moment estimate is $\E[pT^L(h_2f_b)]^4\le e^{-15d/4}$ on the full domain $0\le d\le1/3$.

\paragraph{R35.2 Original finite notation}
Use the original-kernel notation and bounds of
\eqref{eq:odd-common:inputs} and Lemma~\ref{lem:odd-common:spectral}:
\[
 B=-K,\quad u=-g,\quad X=-B^2,\quad
 \tau=d-\delta_2,\quad\theta=13/32,\quad\kappa=19/32.
\]
Put $S_n^-=\one+G_n(B)u$. Its positive convolution is
\eqref{eq:Q188:conv}. The original $X$-spectral measure at $u$
is denoted by $\nu_{\rm sp}$, and its coordinate formulas are
\eqref{eq:odd-common:coordinates}. In particular a constant-vector
component retains its cross term with the same spectral mean.
The genuine message means obey $\E f_n=\E\bar f_n\le1$;
the other two arm messages can have larger means, since
$\E h_2=\E W_2=1+d$.
For the undecorated terminal write
$F_b=h_2f_b$, $H_{b,L}=pT^LF_b$ and
$I_q(Y)=\E[q^{-1/3}|Y|^{8/3}]$.

\paragraph{R35.3 Signed-mean payments and bare paths}
The common signed-mean bound gives
\[
 \E e^{4BV}\le e^{4\E BV+(17/4)\|V\|^2}.\tag{R35.A9}
\]
Corollary~\ref{cor:unit-two-step:bare-paths} gives
\[
 E(p\bar f_2)^4\le e^{d/4},\tag{R35.A10}
\]
\[
 E(p\bar f_n)^4\le1\quad(n\ge3),\qquad pq\le1.\tag{R35.A11}
\]
The $n=3$ proof retains $\|G_2u\|^2=d+\delta_2$.

\paragraph{R35.4 Positive square transport and signed constant-vector terms}

Suppose the actual F satisfies \(0\le F\le (m+R/2)^2,\allowbreak  m>0\) and \(2m+R\ge 0.\) Set

\[
\begin{gathered}
k=1/m,\quad V=ER^2,\quad D_i=EB^iR,\quad a_1=-D_1=EKR,\quad a=ka_1,\\
\quad\delta=kD_0+k^2V/2,\quad\epsilon=\delta-a.
\end{gathered}\tag{R35.A12}
\]

Use the S row of Table~\ref{tab:unit-two-step:parameters}, with
$w=1$, $r=1$, $s_0=m^2$, $s_1=m$, $s_2=1/4$ and $\lambda=1$.
The common finite continuation gives

\[
\begin{gathered}
0\le T^LF/m^2\le Y_L\\
=\bar f_L+kB^LR+\delta \bar f_{L-1}-a \bar f_{L-2}\\
+k\sum_{j=2}^{L-1}D_j\bar f_{L-1-j}\quad(L\ge2).
\end{gathered}\tag{R35.A13}
\]

This upper bound is nonnegative as the result of the positive-square comparison, despite its signed expansion.

For L1, \(TF/m^2\le q+kBR+\delta.\) The sum with p is 2+delta+kBR, yielding

\[
\log E(pTF)^4\le8\log m+4\epsilon+(17/4)k^2V. \tag{R35.A14}
\]

Suppose $R=[a_f(X)-b_f(X)B]u$ and $\rho=\int x b_f$. For $L=2$,

\[
p+Y_2=2+\epsilon+BV_2,\qquad V_2=\delta\one+u+kBR,
\]
\[
EBV_2=-d-k\rho,\quad
\|V_2\|^2=d+2k\rho+k^2\|BR\|^2+\delta^2-2\delta a.
\]

Indeed $\langle u,BR\rangle=\rho$ and $\E BR=-a_1$. The last contribution is exactly \(\epsilon^2-a^2,\allowbreak \) not \(\delta^2\) alone. Dropping only -\(a^2\) gives

\[
\begin{gathered}
\log E(pT^2F)^4\le8\log m+d/4+4\epsilon+(9/2)k\rho\\
+(17/4)k^2\|BR\|^2+(17/4)\epsilon^2.
\end{gathered}\tag{R35.A15}
\]

This is the canonical R calculation, without reversing the terminal pair.

For L3, direct expansion gives

\[
p+Y_3=2+\epsilon-d-k\rho+BV_3,\qquad
V_3=\epsilon\one+(1+\delta)u+Bu+kB^2R,
\]
\[
EBV_3=-(1+\delta)d+kD_3.
\]

Therefore

\[
\begin{gathered}
\log E(pT^3F)^4\le8\log m-8d+4\epsilon-4\delta d\\
+4k(D_3-\rho)+(17/4)\|V_3\|^2.
\end{gathered}\tag{R35.A16}
\]

The nonconstant part of V3 has mean -d-k rho, so its cross term with epsilon1 is -2epsilon(d+k rho), not zero.

For the first connector, Corollary~\ref{cor:unit-two-step:first-square}
with $Z=R$ gives directly
\[
 pTF\le\{m+D_0/2-KR/2+V/(4m)\}^2.\tag{R35.A17}
\]
The square base is nonnegative.

\paragraph{R35.5 Actual low-domain envelopes and whole-frequency bounds}

\paragraph{R35.5.1 \(b=4\)}

Let \(m=1-d\) and \(R=K^2g+K^3g-dg.\) Actual expansion gives \(h_{2}+f_{4}=2m+R\) and

\[
R=(d+X-XB)u,\quad a_f=d+x,\quad b_f=x,
\]
\[
\begin{gathered}
D_0=\delta_2,\quad a_1=d^2+\delta_2,\quad\rho=\delta_3,\quad\\
V=d^3+2d\delta_2+\delta_3+\delta_4,\quad\\
\|BR\|^2=d^2\delta_2+2d\delta_3+\delta_4+\delta_5.
\end{gathered}\tag{R35.A18}
\]

For $d\le1/4$, $\theta^2(1+\theta)=7605/32768<1/4$ gives $\sqrt{a_f^2+xb_f^2}\le u_0=1/2+d$ and $ku_0\le1$. The padded $z_0=\theta+11d/16$ is valid; $b_f\le\theta$ is also available. Exactly,

\[
\begin{gathered}
\epsilon=-kd^2+k^2V/2,\\
|\epsilon|\le d/6,\\
\delta\le25d/24.
\end{gathered}\tag{R35.A19}
\]

For epsilon's upper bound,
\(m^2/6+dm-u0^2/2=(1-4d)(1+8d)/24\ge 0.\)
For its lower bound, \(V\ge 3d^3\); put \(z=d/m\le 1/3\) and use
\(\epsilon/d\ge -z+3z^2/2\ge -1/6.\)
Finally \(\delta/d\le \theta/m+u0^2/(2m^2),\allowbreak \) increasing to \(25/24.\)

\paragraph{R35.5.2 Every finite $b\ge5$, $d\le1/5$}
Use Lemma~\ref{lem:unit-retained-pair} with $n=b$, $\sigma=9/4$,
and $m=1-9d/8$. Denote its retained coefficients by $c_i$, put
$C=c_1+c_2+c_3$, $D=c_2+c_3$, $E=c_3$, and write $R=U$.
Its first actual arm sum gives
\[
 h_2+f_b\le2m+R,\qquad h_2f_b\le(m+R/2)^2.
\]
At $b=5$ the coordinates of $R=(a_f-b_fB)u$ are
$a_f=5d/4+\delta_2+x-x^2$, $b_f=d/4+\delta_2+x$.
For $b\ge6$, write $N=b-2$ and $G_N=a_N+B b_N$. Then
\[
 \begin{gathered}
 a_f=C+x(a_N-E),\qquad b_f=D+xb_N,\\
 a_f-(1-x)b_f\ge c_1+c_2x\ge\tau.
 \end{gathered}\tag{R35.A21}
\]
These are exactly the retained coordinates of R33.8.1--R33.8.3,
whose continuous norm proof keeps the same global $d,\delta_2,E$.
It therefore gives
\[
 \sqrt{a_f^2+xb_f^2}\le u_0=3/8+9d/5,\qquad
 b_f\le z_0=\theta+13d/16.\tag{R35.A23}
\]
The finite propagation uses the seed $+kU$; the signed moments below
are computed for $R=U$, before normalization.

The exact signed spectral moments are

\[
\begin{gathered}
V=\int(a_f^2+xb_f^2),\quad D_0=\int b_f,\quad a_1=\int a_f,\quad\\
\rho=\int xb_f,\quad D_3=\int xa_f,
\end{gathered}
\]
\[
\begin{gathered}
D_{4j+1}=-\int x^{2j}a_f,\quad D_{4j+2}=-\int x^{2j+1}b_f,\quad\\
D_{4j+3}=\int x^{2j+1}a_f,\quad D_{4j+4}=\int x^{2j+2}b_f.
\end{gathered}\tag{R35.A24}
\]

From (R35.A21), \(a1\ge \kappa\) D0+tau d. The finite difference
\(a_f-b_f=c1-Ex+x(a_N-b_N)\) gives

\[
a_1-D_0\le d^2+\theta^3d,\qquad D_3-\rho\le d\delta_2+\delta_5. \tag{R35.A25}
\]

The first possible positive end term is \(x^3,\allowbreak \) or \(x^4\) after the extra x. At b5 the first integrated difference is exactly \(d^2-d3.\)

The complete parameter bounds are

\[
\begin{gathered}
|\epsilon|\le3d/5,\quad-\epsilon<7d/20,\quad\delta\le6d/5,\quad\\
ku_0\le147/155<1,\quad k\rho<3d/10,\quad k^2\theta u_0^2<3/8.
\end{gathered}\tag{R35.A26}
\]

For delta, the target difference after multiplication by \(2m^2\) is
\(463/320-(955/128)d+(2601/1600)d^2,\allowbreak \) decreasing on [0, \(1/5\)], with endpoint \(1577/80000.\)
For epsilon use \(a1\ge \kappa\) D0+tau d. The difference is
\(1867/2560-(64539/20480)d-(118503/51200)d^2,\allowbreak \) decreasing with endpoint \(16519/2560000.\)
The other exact endpoint bounds are
\(k(d+\theta^3)\le 43753/126976<7/20,\allowbreak \)
k theta \(z0\le 1183/3968<3/10,\allowbreak \) and
\(k^2\theta u0^2\le 280917/768800<3/8.\)
All belong to the same actual R.

\paragraph{R35.5.3 L3 frequency inequality and its certificate}

In both low domains,

\[
(1+\delta-kxa_f)^2+x(1+kxb_f)^2\le(1+\delta)^2+x. \tag{R35.A27}
\]

Use \(a_f^2+x b_f^2\le u0 a_f+x\) z0 \(b_f\) and \(a_f\ge (1-x)b_f.\) Since 2-kx \(u0>0,\allowbreak \) the unshifted difference is at most

\[
kx b_f[-2+4x+kx\{u_0(1-x)+xz_0\}].
\]

The quantity x[u0(1-x)+xz0] increases with x, since its derivative is \ensuremath{\ge}\(u0(1-2x)>0,\allowbreak \) and the parameters increase with d. At the upper corners, the negative-bracket margins are

\[
323/49152>0\quad(b4,\ z_0=\theta+11d/16),\qquad
79745/3174400>0\quad(b\ge5).
\]

For $b=4$, the sharper input $b_f=x\le\theta$ gives margin $1091/24576$. With the padded input $z_0=\theta+11d/16$, the valid margin is $323/49152$. Both are positive. The scalar shift contributes $-2\delta kx a_f\le0$.

The actual constant coordinate consequently gives

\[
\|V_3\|^2\le(1+\delta)^2d+\delta_2+\epsilon^2-2\epsilon(d+k\rho). \tag{R35.A28}
\]



\paragraph{R35.6 All low-domain short connectors}\label{part:R35.6}

For b4, \(d\le 1/4,\allowbreak \) (R35.A14) and exact epsilon yield

\[
\log EH_{4,1}^4/d\le-8-8d+\frac{25}{4}\left(\frac{1/2+d}{1-d}\right)^2.\tag{R35.A29}
\]

Here \(8\log(1-d)\le -8d-4d^2\) and \(k\ge 1.\) The right side is convex, with endpoint values -\(103/16,\allowbreak -15/4.\) For $b\ge5$, the complete-host $L=1$ conclusion is Proposition~\ref{prop:unit-retained:whole-L1}.

For canonical L2, \(b\ge 5,\allowbreak \) use \(\lVert BR\rVert ^2\le \theta\) V in (R35.A15), keeping the negative quadratic logarithm term:

\[
\begin{gathered}
\log EH_{b,2}^4/d\\
\le-127/20-(1413/400)d+(9/2)\theta z_0/m\\
+(17/4)\theta(u_0/m)^2.
\end{gathered}\tag{R35.A31}
\]

This is convex. At zero the two positive terms sum to less than one, giving below -\(107/20.\) At \(d1/5,\allowbreak \) (R35.A26) gives below -\(16451/4000<-15/4.\) The coefficient -\(1413/400\) is exactly -\(81/16+(17/4)(9/25),\allowbreak \) retaining the signed constant-vector budget.

For b4, L2,

\[
\|BR\|^2/d\le\theta/16+\theta^2/2+\theta^3+\theta^4
=212017/1048576<13/64.
\]

Use \(k\le 4/3,\allowbreak  \epsilon^2/d\le 1/144\) and \(\rho/d\le \theta^2<1/6\) in (R35.A15):

\[
\log EH_{4,2}^4/d
\le-85/12+1+221/144+17/576=-2603/576<-15/4. \tag{R35.A32}
\]



For b4, L3, the scalar shift plus the mean in (R35.A16) simplifies exactly to
-\(2d-kd^2+kV/2.\) By (R35.A19), k \(\rho<2d/9\) and (R35.A28),

\[
\|V_3\|^2/d\le(1+25d/24)^2+\theta+(47/108)d
\le58163/27648<17/8.
\]

The coefficient \(47/108=1/36+11/27\) pays \(\epsilon^2\) and its possible positive cross term. Also \(2kV/d\le 3/2,\allowbreak \) so

\[
\log EH_{4,3}^4/d\le-16+3/2+(17/4)(17/8)=-175/32. \tag{R35.A33}
\]

For \(b\ge 5,\allowbreak \) L3, if \(\epsilon\ge 0\) its cross term in (R35.A28) is nonpositive; if \(\epsilon<0\) it is <\(91d^2/100.\) Hence

\[
\|V_3\|^2/d\le1+\theta+(367/100)d+(36/25)d^2<9/4.
\]

Also \(4k(D3-\rho)/d\le 4k(\theta d+\theta^4)<3/5,\allowbreak \) since
4k \(\theta<17/8,\allowbreak \) 4k \(\theta^4<3/20\) and \(d\le 1/5.\) Therefore (R35.A16) is at most

\[
[-9-8+12/5+3/5+(17/4)(9/4)]d=-71d/16. \tag{R35.A34}
\]

This completes the low-domain $L=2,3$ ordinary-fourth rows with $\E H^4\le e^{-15d/4}$, together with the $b=4,L=1$ row. For $b\ge5,L=1$, Proposition~\ref{prop:unit-retained:whole-L1} provides the required whole-domain outlet.

\paragraph{R35.7 Every finite \(L\ge 4,\allowbreak \) including the real constant coordinate}

Rewrite (R35.A13) using the positive convolution \eqref{eq:Q188:conv}:

\[
\begin{gathered}
Z_L^*=S_L^-+\delta S_{L-1}^- -aS_{L-2}^-+kB^LR\\
=\nu_0+Q_L(B)u,\quad\nu_0=1+\epsilon.
\end{gathered}
\]

Its difference from \(Y_L\) is exactly \(\sum_{j=3}^L\gamma_j\bar f_{L-j}\), where

\[
\begin{gathered}
\gamma_3=d-kD_2,\quad\gamma_4=d+\delta d-kD_3,\quad\\
\gamma_j=\beta_j+\delta\beta_{j-1}-a\beta_{j-2}-kD_{j-1}\ (j\ge5).
\end{gathered}\tag{R35.A35}
\]

Taking means gives \(\alpha_L:=nu0-EZ_L^*= \gamma_{L+1}\) for \(L\ge 4.\) This uses \(ES_n^-=1- \beta_{n+1}\), not a mean-zero assumption.

The exact finite polynomials are

\[
\begin{gathered}
P_n=G_n+kB^n(a_f-b_fB),\qquad\\
Q_n=P_n+\epsilon G_{n-2}+\delta B^{n-2}\quad(n\ge3).
\end{gathered}\tag{R35.A36}
\]

For every such n, \(\lVert P_n\rVert \le 1.\) At n3 the coordinates are
\(1-x-kx^2b_f,\allowbreak \) 1-kx \(a_f,\allowbreak \) nonnegative and no larger than the G3 coordinates, because \(ku0\le 1\) and \(kz0\le 37/48\) for b4 or \(91/124\) for \(b\ge 5\); \(1-x-x^2>0.\)
At n4,
|\(P4|\le (1-x)(1+x/2)+x^2\le 1.\)
At n5,
|\(G_{5}|\le 1-x/2+x^2/2\) and \(x^{5/2}\)\ensuremath{\le}\(13x^2/20,\allowbreak \) giving
|\(P5|\le 1-x/2+23x^2/20\le 1\) since \(\theta<10/23.\)
The squared comparison for G5 has exact difference \(x^2(1+2x-3x^2)/4\ge 0.\)
For every \(n\ge 6,\allowbreak \)

\[
|P_n|\le(1+x^3)/\sqrt{1+x}+x^3\le1.
\]

For \(x>0\) the last condition is equivalent to \(x(1+\sqrt{1+x})\le 1,\allowbreak \) which follows from its bound (\(13/32)(11/5)=143/160.\) The zero frequency is direct.

Using (R35.A19),(R35.A26), \(\sqrt{\theta}<13/20\) and \(\lVert G_{2}\rVert <6/5\) gives

\[
\begin{gathered}
\|Q_n\|\le1+d\quad(b4),\qquad\\
\|Q_n\|\le1+7d/5\quad(b\ge5),\qquad n\ge3.
\end{gathered}\tag{R35.A37}
\]

At n3 the added costs are at most 27d/32 and 69d/50. At n4 use (\(6/5)|\epsilon|+\theta\) delta; at \(n\ge 5\) use |epsilon|+\(\theta^{3/2}\)delta. Both are smaller than the stated respective budgets. These bounds cover the finite-length tails.

Let \(h=1/2-6d/5>0.\) Every \(\gamma_j\ge h\) d.
\(D2\le 0\) gives \(gamma3\ge d.\) Since \(ka1\le ku0 d\le d\) and \(D3\le \theta\) a1, \(gamma4\ge \kappa\) d.
For \(j\ge 5,\allowbreak \) each positive \(D_{j-1}\) is \ensuremath{\le}\(\theta^2z0\) d. The first positive even class is D4; the first relevant positive odd class is \(D7\le \theta^3u0 d\le \theta^2z0\) d, and later powers decrease. The needed theta \(u0\le z0\) holds throughout both domains.

For b4 the lower coefficient is kappa-ku0 d-k \(\theta^2z0.\) Subtract h and multiply by m; the result is

\[
875/32768+(40369/81920)d-(11/5)d^2,
\]

a concave quadratic with positive endpoints \(875/32768\) and \(4063/327680.\)
For \(b\ge 5\) use \(a\le d.\) The corresponding difference is

\[
875/32768-(3241/81920)d-(9/40)d^2,
\]

decreasing with positive endpoint \(40101/4096000.\) Therefore \(\gamma_j\ge hd\) for every \(j\ge 3\) and \(\alpha_L\ge hd.\)

Retain the last two actual terms \(j=L,\allowbreak \) L-1:

\[
0\le T^LF/m^2\le Z=Z_L^*-hd(1+q).
\]

The finite recurrence \(Q_L=nu0+ BQ_{L-1}\) gives

\[
\begin{gathered}
Z=Cq-hd+N_0,\\
C=\nu_0-hd,\\
N_0=BQ_{L-1}u,\\
EN_0=-\alpha_L\le-hd.
\end{gathered}\tag{R35.A38}
\]

This displays the real constant C. Once \(C>0,\allowbreak \) weighted AM--GM between actual nonnegative p, Z/C yields

\[
\begin{gathered}
H_{b,L}\le(c_*+N/2)^2,\quad\\
c_*=m\sqrt C-\frac{mhd}{2\sqrt C},\quad N=\frac m{\sqrt C}N_0.
\end{gathered}\tag{R35.A39}
\]

The base equals \(m(\sqrt{C}p+Z/\sqrt{C})/2\ge 0.\) Writing \(N=BV_*,\allowbreak \) (R35.A9) gives fourth exponent

\[
8(c_*-1)+4EN+(17/4)\|V_*\|^2. \tag{R35.A40}
\]

For b4,
\(1-2d/3+6d^2/5\le C\le 1\) and \(C-m^2\ge 4d/3+d^2/5>0\); in particular \(C>0.\)
By (R35.A37), \(\lVert V_*\rVert ^2\le (1-d^2)^2d/C\le 17d/16.\) The full comparison is

\[
(17/16)(1-2d/3+6d^2/5)-(1-d^2)^2
\ge1/16-(17/24)d+(257/80)d^2>0.
\]

Here \(d^4\le d^2/16\) and the discriminant is negative because \(289/576<257/320.\)
Also \(m/\sqrt{C}\ge 3/4\) and \(h\ge 1/5,\allowbreak \) so \(EN\le -3d/20.\)
Using m \(\sqrt{C}\le (m^2+C)/2,\allowbreak \)

\[
c_*\le1-(65/48)d+(31/20)d^2\le1-29d/30.
\]

Thus (R35.A40) is at most
[-\(116/15-3/5+289/64\)]\(d=-733d/192<-15d/4.\)

For \(b\ge 5,\allowbreak \)

\[
C_{\rm low}=1-11d/10+6d^2/5\le C\le1+d/10+6d^2/5\le267/250,
\]
\[
C-m^2\ge(23/20)d-(21/320)d^2>0.
\]

Again \(C>0.\) Expanding,

\[
m^2(1+7d/5)^2=(1+11d/40-63d^2/40)^2
\le1+11d/20-(4919/1600)d^2.
\]

The discarded tail is \(d^3\)[-\(693/800+(3969/1600)d\)]\ensuremath{\le}0. Subtracting this upper bound from (\(5/4)C_low\) leaves
\(1/4-(77/40)d+(7319/1600)d^2>0,\allowbreak \) whose discriminant comparison is \(5929<7319.\)
Thus \(\lVert V_*\rVert ^2\le 5d/4.\) Moreover \(m/\sqrt{C}\ge 149/200,\allowbreak \) since

\[
961/1600-(149/200)^2(267/250)=78583/10000000>0.
\]

With \(h\ge 13/50,\allowbreak  EN\le -(149/200)(13/50)d<-19d/100.\)
Using \(\sqrt{C}\le 1+(C-1)/2\) and dropping only a nonpositive cubic,

\[
c_*\le1-(1009/800)d+(3963/4000)d^2\le1-21d/20.
\]

The endpoint coefficient is \(10631/10000>21/20.\) Hence (R35.A40) is at most
[-\(42/5-19/25+85/16\)]\(d=-1539d/400<-15d/4.\)
Every finite low-domain \(L\ge 4\) is now paid.

\paragraph{R35.8 Actual high-domain terminal means and every \(L\ge 2\)}

\paragraph{R35.8.1 \(b=4,\allowbreak \) including the remaining L1 domain}

Original adjunction and finite convolution give

\[
t_4:=E(h_2f_4)=1-3d+\delta_2-d^2+\delta_3
\le1-(2487/1024)d-d^2.
\]

For \(d\ge 1/4\) this upper bound decreases, with value \(1353/4096<1/3\) at \(1/4.\) Thus \(TF4\le 2/3.\) Positivity gives \(H_{4,\allowbreak  2}\)\ensuremath{\le}(\(2/3)pq,\allowbreak  H_{4,\allowbreak  3}\)\ensuremath{\le}(\(2/3)p\bar f_2\) and \(H_{4,\allowbreak  L}\)\ensuremath{\le}(\(2/3) p\bar f_{L-1}\) for \(L\ge 4.\)
Equations (R35.A10)-(R35.A11), \(4\log(2/3)\le -4/3\) and \(d\le 1/3\) prove all these ordinary-fourth bounds \ensuremath{\le}exp(-15d/4). In the only cost-bearing row L3, -\(4/3+d/4\le -15d/4\) is exactly \(d\le 1/3.\)

For b4, L1 use (R35.A17) with exact \(D0=\delta_2\) and \(EZ=d^2+\delta_2.\) Applying (R35.A9) to -R cancels the linear d2 terms and gives

\[
\log EH_{4,1}^4\le-8d-4d^2+\{2/(1-d)+17/4\}V. \tag{R35.A41}
\]

The bound \(V/d\le d^2+2\theta d+\theta^2+\theta^3<5/8\) holds throughout \(d\le 1/3:\) its left side increases, and at \(1/3\) the positive cross difference is
\(35*32768-7605*144=51760.\)
For \(d\ge 1/4\) the coefficient in (R35.A41) is \ensuremath{\le}\(29/4,\allowbreak \) giving
\(\log(\E H_{4,\allowbreak 1}^4)/d\le-9+145/32=-143/32<-15/4\).
Combined with \mbox{\hyperref[part:R35.6]{part R35.6}} this proves the b4, L1 full-domain row.

\paragraph{R35.8.2 Every finite \(b\ge 5\) on \(1/5\le d\le 1/3\)}

Let \(t_b=E(h_{2} f_b),\allowbreak \) an actual nonnegative mean. Adjunction gives
\(t_b=E\)[p \(f_{b+1}\)]=2\(\mu_{b+1}\)-\(\mu_{b+2}\).
Multiply the original convolution by \(h_{2}=1-g-Kg\) to obtain

\[
\begin{gathered}
t_b=1-\beta_{b+1}-\beta_{b+2}+\beta_{b+3}\\
-\sum_{j=3}^b\beta_jt_{b-j}.
\end{gathered}\tag{R35.A42}
\]

In particular actual \(t0=1+d\) is retained. Direct finite expansion gives

\[
\begin{aligned}
t_0&=1+d,&t_1&=1,&t_2&=1-d-\delta_2,\\
t_3&=1-2d-d^2,\\
t_4&=1-3d+\delta_2-d^2+\delta_3,\\
t_5&=1-4d+2(1+d)\delta_2,\\
t_6&=1-5d+2d^2+d^3+(3+2d)\delta_2-\delta_3-\delta_4,\\
t_7&=1-6d+5d^2+2d^3+4\delta_2-\delta_2^2-2(1+d)\delta_3,\\
t_8&=1-7d+9d^2+2d^3+(5-4d-3d^2)\delta_2-\delta_2^2
\\ &\quad -(3+2d)\delta_3+\delta_4+\delta_5.
\end{aligned} \tag{R35.A43}
\]

For explicit finite-end checks, b7 has base 1-d+d2-d3 and loss

\[
d(t_4+t_3)+\tau(t_2+t_1+t_0)+\delta_3t_0
=5d-5d^2-2d^3-3\delta_2+\delta_2^2+\delta_3+2d\delta_3.
\]

At b8 the base is 1-d+d2-d3+d4+d5 and loss is

\[
\begin{gathered}
d(t_5+t_4)+\tau(t_3+t_2+t_1+t_0)+\delta_3(t_1+t_0)\\
=6d-9d^2-2d^3+(4d+3d^2-4)\delta_2+\delta_2^2+2\delta_3+2d\delta_3.
\end{gathered}
\]

The four short cases have the following upper polynomials for \(t_b-(1-d)/2:\)


\begingroup\small
\setlength{\tabcolsep}{3pt}
\setlength{\LTleft}{0pt}\setlength{\LTright}{\fill}
\begin{longtable}{@{}p{0.1200\linewidth}p{0.8000\linewidth}@{}}
\caption{The terminal pairs $(2,b)$ for every $b\ge4$ at a unit root: parameter values.}\label{tab:parameters:r2-c035:2}\\
\toprule
Entry & Upper polynomial \\
\midrule\endfirsthead
\toprule
Entry & Upper polynomial \\
\midrule\endhead
\bottomrule\endfoot
5 & \(1/2-(43/16)d+(13/16)d^2\); Complete continuous payment: decreasing; left value -\(1/200\). \\
\addlinespace[2pt]
6 & \(1/2-(105/32)d+(45/16)d^2+d^3\); Complete continuous payment: derivative \ensuremath{\le}-\(103/96\); left value -\(143/4000\). \\
\addlinespace[2pt]
7 & \(1/2-(31/8)d+5d^2+2d^3\); Complete continuous payment: convex; endpoints -\(59/1000,\allowbreak -35/216\). \\
\addlinespace[2pt]
8 & \(1/2-(143/32)d+(59/8)d^2+(25/32)d^3\); Complete continuous payment: convex; endpoints -\(37/400,\allowbreak -61/432\). \\
\end{longtable}
\endgroup


For b8, first group \(d4+d5\le (\theta+\theta^2)d3\le (3+2d)d3,\allowbreak \) then use the positive coefficient \(5-4d-3d^2\) with \(\delta_2\le \theta\) d. At b7 the omitted terms are explicitly nonpositive.

For every finite \(b\ge 9,\allowbreak \) the seven actual coefficients of t0,..., t6 in (R35.A42) have indices at least three. The base is \ensuremath{\le}1-tau+d3, from the finite beta increments. Their exact sum is

\[
\sum_{j=0}^6t_j=7-14d+d^3+(5+4d)\delta_2-\delta_4.
\]

Thus

\[
t_b\le1-\tau[8-14d+d^3+(5+4d)\delta_2-\delta_4]+\delta_3. \tag{R35.A44}
\]

Let P be the amount subtracted from one. For the original global \(v=\delta_2/d\) in [d, theta],

\[
P\ge d[(1-v)(A+Cvd)-\theta v],\qquad
A=8-14d+d^3,\quad C=5+4d-\theta^2.
\]

The bracket decreases in v because its derivative is
-\(A+Cd(1-2v)-\theta<0,\allowbreak \) using \(A\ge 10/3\) and \(Cd\le 19/9.\)
Take the legitimate global endpoint \(v=\theta,\allowbreak \) and discard only nonnegative \(d^3,\allowbreak  d^4\) terms. Then

\[
\begin{gathered}
P-(1+d)/2\ge-1/2+(4183/1024)d-(7493391/1048576)d^2\\
\ge-1/2+4d-(29/4)d^2>0.
\end{gathered}\tag{R35.A45}
\]

The last concave polynomial has endpoints \(1/100,\allowbreak  1/36.\) Hence every finite \(b\ge 5\) satisfies

\[
t_b\le(1-d)/2,\qquad 0\le TF_b\le1-d. \tag{R35.A46}
\]

All finite b cases are exhausted.

For \(L\ge 2,\allowbreak \) positivity gives \(H_{b,\allowbreak  L}\)\ensuremath{\le}(1-d)\(p\bar f_{L-1}\). Use \(pq\le 1\) at L2, (R35.A10) at L3, and (R35.A11) at \(L\ge 4.\) Since \(4\log(1-d)\le -4d,\allowbreak \) the exponents are respectively -4d,-15d/4,-4d. Combined with parts R35.6-7 and b4 high,

\[
EH_{b,L}^4\le e^{-15d/4}\qquad(b\ge4,\ L\ge2,\ 0\le d\le1/3). \tag{R35.A47}
\]

\paragraph{R35.9 Every $b\ge5$, $L=1$, on the complete host domain}
Apply Proposition~\ref{prop:unit-retained:whole-L1} to the actual
$H_{b,1}=pT(h_2f_b)$. It proves
\[
 I_q(T^2H_{b,1})^3\le e^{-d/500}\qquad(0\le d\le1/3).
\]
The same $\sigma=2$ retained vector and the same canonical direction
are used throughout this interval. Its $b=5$ and $b\ge6$ data are
the two integer cases in Lemma~\ref{lem:unit-retained:whole-data}.

\paragraph{R35.11 Exact C094 bridge, curvature and logarithm payment}

\paragraph{R35.11.1 The original outer output}
Use Corollary~\ref{cor:unit-two-step:output} and
paragraph~\ref{part:R29.11}. On the original measure, every real
$Y$ satisfies
\[
 I_q(T^2Y)^3\le B_2(E|Y|^4)^2,\tag{R35.A61}
\]
\[
 B_2=[1-2d^2+(8-2d)\delta_2](1+1024\delta_2/855)^8
       \le e^{c_2d-5d^2},\qquad c_2=24427/3420.\tag{R35.A62}
\]

\paragraph{R35.11.2 Completion of the ordinary-fourth rows}
Every row with $\E H^4\le e^{-15d/4}$ yields
\[
 \log I_q(T^2H)^3\le(c_2-15/2)d-5d^2
 =-(1223/3420)d-5d^2\le-d/500.\tag{R35.A63}
\]
Equation (R35.A47) covers all $L\ge2$ and $b\ge4$.
Paragraph R35.8.1 and the low-domain bound in R35.6 cover $b=4,L=1$.
For $b\ge5,L=1$ use Proposition~\ref{prop:unit-retained:whole-L1}
on its complete host domain. Hence
\begin{equation}\tag{R35.A68}
 I_q(T^2[pT^L(h_2f_b)])^3\le e^{-d/500}\le1
 \qquad(b\ge4,\ L\ge1).
\end{equation}

\paragraph{R35.12 Original-tree realization, correlated KL, full reversal and all maps}

The direction is fixed before (mu, K):
the terminal b-arm points outward;
the two-arm has first edge toward the terminal and second toward its leaf;
the internal root's unit arm points outward;
all L inner edges point from terminal to internal root;
both outer edges point from internal root to receiving root.
Every genuine equal pair uses \(f_l,\allowbreak  h_l\) at its original vertex.
Hence the exact original output is

\[
H_D=E_rpT^L(E_t h_2f_b),\qquad U_D=T^2H_D.
\]

No analytic case changes this direction, including L2.
Since \(0<E_t,\allowbreak  E_r\le 1\) on the positive support and T is positive,

\[
0<H_D\le pT^L(h_2f_b),\qquad
0<U_D\le T^2[pT^L(h_2f_b)].
\]

The positive \(I_q\) weight makes (R35.A68) apply to the actual decorated output.
The multiplier is the specified product of equal-pair factors.
Without decorations the component has exactly L+b+5 original edges:
terminal 2+b, inner L, root side 1, outer 2.
Each equal pair adds its actual 2l edges and original interior vertices.
All subdivision vertices and independently specified core lengths remain in the integral.
The local outer length is two; other independently specified outer lengths retain their original values.

The outer kernel remains

\[
T^2(x,y)=E_z[(1-K(x,z))(1-K(z,y))].
\]

Adjoints and Holder use the original mu, \(T*=A.\)
Paragraph~\ref{part:R29.11} controls the inverse weight through the
same-kernel centered-row bound as $q$ becomes small.
Constants are independent of min mu; q may be arbitrarily small across finite models.

At \(d=0,\allowbreak  K1=0\) on the positive support but K need not be zero.
\(A1=T1=1,\allowbreak \) so leaf integration gives every actual original tree message one.
Then \(H_D=U_D=q=1\) and the claimed result is equality.
All divisions by d were confined to \(d>0.\)
Scalar zero-variance cases and frequency \(x=0\) were also included.

Complete reversal is substitution \(K\mapsto\) -K in the same fixed structure.
The terminal becomes W2 \(\bar f_b\); inner \(T^L\) becomes \(A^L\); unit p becomes q;
all equal pairs become \(\bar f_l W_l\); outer \(T^2\) becomes \(A^2\); \(I_q\) becomes \(I_p.\)
The same loss follows because d is unchanged.
Swapping only one weight or one connector while retaining other orientations is not justified.
